\section{The code}
\label{sec:code}

The library is in \cmd{src/} (C11, Linux kernel coding style, no global
state: one design point per thread), its headers in
\cmd{include/mapping/}, the programs in \cmd{tools/}.

\begin{center}\small\row
\begin{tabularx}{\textwidth}{@{}>{\raggedright\arraybackslash}p{30mm}X@{}}
\toprule\rowcolor{white}
file & role (paper section) \\
\midrule
\cmd{model.c} & tasks, chains, cores; WCET on a core, with contention (System model) \\
\cmd{context.c} & the parameters of a design point, the platform and its ECUs \\
\cmd{json.c}, \cmd{io.c} & the JSON reader; the application and configuration files; the templates \\
\cmd{generator.c} & synthetic applications, replicas, the detection chain \\
\cmd{partition.c} & cutting into sub-chains (Step 1) \\
\cmd{let.c} & the LET task of a core and the core loads \\
\cmd{assign.c} & initial cores, allocation rules, memory reservation, contention update (Step 2) \\
\cmd{placement.c} & consolidation, relocation moves and the optional swap moves (Steps 3 and 5) \\
\cmd{sched.c} & response times, priority assignment, scaling factor (Timing analysis) \\
\cmd{repair.c} & deadline-driven local search (Step 6) \\
\cmd{communication.c} & segments, communications, end-to-end bounds, incremental analysis \\
\cmd{network.c} & topologies, routes, CBS/TAS/ATS service, FRER, fixed point \\
\cmd{stats.c} & metrics, verification, lower bound \\
\cmd{pipeline.c} & one design point, contention fixed point, sizing \\
\cmd{lifecycle.c} & the SOP and the software updates (Life of a platform) \\
\cmd{search.c} & simulated annealing and genetic algorithm (baselines) \\
\cmd{experiment.c} & parallel experiments and CSV output \\
\bottomrule
\end{tabularx}
\end{center}

\paragraph{Adding an experiment.} In \cmd{tools/experiment.c}, write a
function that fills a list of points with \cmd{add\_point()} (one call per
point and rule, setting the parameters to vary; \cmd{mp\_params\_template()}
gives a template platform), and add a line to the \cmd{experiments[]}
table with its name, description, default number of instances and the
columns of its summary.

\paragraph{Adding a key to the files.} Add the field to \cmd{struct
mp\_params} (\cmd{context.h}) with its default in \cmd{mp\_params\_default()},
read it in the right section of \cmd{io.c} (with its key in the list of
known keys and its range), write it in \cmd{mp\_write\_configuration()},
and document it in this manual. The \cmd{roundtrip} and \cmd{unit\_io}
tests then check that what is written is read back.
